\chapter{Graphs And Subgraphs}
%%%%%%%%%%%%%%%%%%%%%%%%%%%%%%%%%%%%%%%%%%%%%%%%%%%%%%%%%%%%%%%%%%%%%%%%%%%%%%%
\section{Graphs And Subgraphs}

% Generated by scripts/blueprint_restate.py from TCSlib.GraphTheory.Core.GraphsAndSubgraphs.
% Each body is the STATEMENT only, taken from the declaration's Lean docstring:
% the one-line summary plus the verbatim book statement.  Proof, proof plan,
% significance commentary and formalisation notes are excluded by construction.

\begin{theorem}[Adjacency-preserving bijection characterises graph isomorphism]
\lean{SimpleGraph.nonempty_iso_iff}
\label{SimpleGraph.nonempty_iso_iff}
\leanok
Let $G$ be a simple graph on a finite vertex set $V$ and $H$ a simple graph on a vertex
set $W$. Then $G$ and $H$ are isomorphic if and only if there is a bijection $\theta
\colon V \to W$ such that for all vertices $u, v \in V$, the pair $u, v$ is adjacent in
$G$ if and only if $\theta(u), \theta(v)$ is adjacent in $H$.
\end{theorem}

\begin{theorem}[Maximal edge count characterises the complete graph]
\lean{SimpleGraph.edgeCard_eq_choose_two_iff_top}
\label{SimpleGraph.edgeCard_eq_choose_two_iff_top}
\leanok
Let $G$ be a simple graph on a finite vertex set $V$ with $n = \abs{V}$ vertices. Then
$G$ has exactly $\binom{n}{2}$ edges if and only if $G$ is the complete graph on $V$,
that is, every pair of distinct vertices is adjacent.
\end{theorem}

\begin{theorem}[Edge count of a complete bipartite graph]
\lean{SimpleGraph.completeBipartite_edgeCard}
\label{SimpleGraph.completeBipartite_edgeCard}
\leanok
For natural numbers $m$ and $n$, let $K_{m,n}$ be the complete bipartite graph whose two
parts have $m$ and $n$ vertices, with an edge joining every vertex of the first part to
every vertex of the second and no edges within either part. Then $K_{m,n}$ has exactly
$mn$ edges.
\end{theorem}

\begin{theorem}[Edge bound for bipartite graphs]
\lean{SimpleGraph.bipartite_edgeCard_le}
\label{SimpleGraph.bipartite_edgeCard_le}
\leanok
Let $G$ be a simple graph on a finite vertex set $V$ that is bipartite, i.e. admits a
proper $2$-colouring. Then the number of edges of $G$ satisfies \[ 4\,\abs{E(G)} \le
\abs{V}^{2}. \]
\end{theorem}

\begin{definition}[Ex 1.2.10: the $k$-cube --- vertices are $k$-bit strings, adjacent iff they differ in exactly one coordinate]
\lean{SimpleGraph.hypercube}
\label{SimpleGraph.hypercube}
\leanok
Ex 1.2.10: the \texttt{k}-cube --- vertices are \texttt{k}-bit strings, adjacent iff
they differ
in exactly one coordinate.

\emph{The \texttt{k}-cube is the graph whose vertices
are the ordered \texttt{k}-tuples of \texttt{0}s and \texttt{1}s, two vertices being
joined if and
only if they differ in exactly one coordinate.}
\end{definition}

\begin{theorem}[Edge count of the k-cube]
\lean{SimpleGraph.hypercube_edgeCard}
\label{SimpleGraph.hypercube_edgeCard}
\leanok
\uses{SimpleGraph.hypercube}
For every $k \in \bbn$, the $k$-cube --- the graph whose vertices are the length-$k$
binary strings, two strings being joined by an edge precisely when they differ in
exactly one coordinate --- has exactly $k \cdot 2^{k-1}$ edges.
\end{theorem}

\begin{theorem}[The k-cube is bipartite]
\lean{SimpleGraph.hypercube_bipartite}
\label{SimpleGraph.hypercube_bipartite}
\leanok
\uses{SimpleGraph.hypercube}
For every $k \in \bbn$, the $k$-cube $Q_k$ --- the graph whose vertices are the
length-$k$ strings of bits, two strings being adjacent precisely when they differ in
exactly one coordinate --- admits a proper colouring with two colours; that is, $Q_k$ is
bipartite.
\end{theorem}

\begin{theorem}[Order of a self-complementary graph modulo four]
\lean{SimpleGraph.selfComplementary_card_mod_four}
\label{SimpleGraph.selfComplementary_card_mod_four}
\leanok
Let $G$ be a self-complementary simple graph on a finite vertex set $V$; that is, $G$ is
isomorphic to its complement $\overline{G}$. Then the number of vertices $\abs{V}$
satisfies $\abs{V} \equiv 0 \pmod 4$ or $\abs{V} \equiv 1 \pmod 4$.
\end{theorem}

\begin{theorem}[Induced subgraph of a complete graph]
\lean{SimpleGraph.induce_completeGraph}
\label{SimpleGraph.induce_completeGraph}
\leanok
Let $V$ be a vertex set and let $K_V$ denote the complete graph on $V$, in which every
pair of distinct vertices is adjacent. For any subset $s \subseteq V$, the subgraph of
$K_V$ induced on $s$ is again a complete graph, namely the complete graph $K_s$ on the
vertex set $s$: any two distinct vertices of $s$ are adjacent in the induced subgraph
precisely because they are adjacent in $K_V$.
\end{theorem}

\begin{theorem}[Subgraphs inherit 2-colourability]
\lean{SimpleGraph.subgraph_bipartite}
\label{SimpleGraph.subgraph_bipartite}
\leanok
Let $G$ and $H$ be simple graphs on a common finite vertex set, and suppose $H$ is a
subgraph of $G$ (every edge of $H$ is an edge of $G$). If $G$ admits a proper vertex
colouring with two colours, then so does $H$.
\end{theorem}

\begin{theorem}[Two-colourability passes to induced subgraphs]
\lean{SimpleGraph.induce_bipartite}
\label{SimpleGraph.induce_bipartite}
\leanok
Let $G$ be a simple graph on a finite vertex set $V$ that is $2$-colourable, meaning it
admits a proper colouring with two colours. Then for every subset $s \subseteq V$, the
induced subgraph $G[s]$ is again $2$-colourable.
\end{theorem}

\begin{theorem}[Bipartiteness passes to induced subgraphs]
\lean{SimpleGraph.subgraph_induce_bipartite}
\label{SimpleGraph.subgraph_induce_bipartite}
\leanok
Let $G$ be a simple graph on a finite vertex set, and suppose $G$ is bipartite, meaning
its vertices can be properly colored with two colors so that adjacent vertices receive
different colors. Let $H$ be any subgraph of $G$ — a simple graph on the same vertex set
every edge of which is also an edge of $G$ — and let $s$ be any set of vertices. Then
the subgraph of $H$ induced on $s$, obtained by keeping only the vertices in $s$
together with those edges of $H$ joining two vertices of $s$, is again bipartite.
\end{theorem}

\begin{theorem}[Degree-sum formula (handshaking lemma)]
\lean{SimpleGraph.sum_degrees_eq_two_mul_edgeCard}
\label{SimpleGraph.sum_degrees_eq_two_mul_edgeCard}
\leanok
Let $G$ be a simple graph on a finite vertex set $V$, and for each vertex $v$ write
$d(v)$ for its degree, the number of vertices adjacent to $v$. Then \[ \sum_{v \in V}
d(v) = 2\,\abs{E(G)}, \] where $\abs{E(G)}$ is the number of edges of $G$.
\end{theorem}

\begin{theorem}[Parity of the odd-degree vertex count]
\lean{SimpleGraph.even_card_odd_degree}
\label{SimpleGraph.even_card_odd_degree}
\leanok
In a finite simple graph $G$, the number of vertices whose degree is odd is even.
\end{theorem}

\begin{theorem}[Minimum and maximum degree bounds on edge count]
\lean{SimpleGraph.degree_bounds}
\label{SimpleGraph.degree_bounds}
\leanok
Let $G$ be a simple graph on a finite, nonempty vertex set $V$, and write $n = |V|$ for
the number of vertices and $m$ for the number of edges. Let $\delta(G)$ and $\Delta(G)$
denote the minimum and maximum degree of $G$, respectively. Then
\[
n\,\delta(G) \;\le\; 2m \;\le\; n\,\Delta(G).
\]
\end{theorem}

\begin{theorem}[Equal parts in a regular bipartite graph]
\lean{SimpleGraph.regular_bipartite_card_eq}
\label{SimpleGraph.regular_bipartite_card_eq}
\leanok
Let $G$ be a finite simple graph, and let $k \ge 1$ be a positive integer such that
every vertex of $G$ has degree exactly $k$. Suppose that $G$ is bipartite with parts $s$
and $t$; that is, $s$ and $t$ are disjoint sets of vertices whose union contains every
vertex, and every edge of $G$ joins a vertex of $s$ to a vertex of $t$. Then the two
parts have the same size, $\abs{s} = \abs{t}$.
\end{theorem}

\begin{theorem}[Bipartite spanning subgraph preserving half of every degree]
\lean{SimpleGraph.exists_bipartite_spanning_subgraph}
\label{SimpleGraph.exists_bipartite_spanning_subgraph}
\leanok
Let $G$ be a finite simple graph with vertex set $V$. Then $G$ has a spanning subgraph
$H$ that is bipartite (that is, $2$-colourable) and in which every vertex keeps at least
half of its original degree: writing $d_G(v)$ for the degree of $v$ in $G$ and $d_H(v)$
for the number of neighbours of $v$ in $H$, one has $d_G(v) \le 2\,d_H(v)$ for every $v
\in V$.
\end{theorem}

\begin{theorem}[Edge count of a line graph]
\lean{SimpleGraph.lineGraph_edgeCard}
\label{SimpleGraph.lineGraph_edgeCard}
\leanok
Let $G$ be a finite simple graph, and let $L(G)$ denote its line graph, whose vertices
are the edges of $G$, with two such vertices adjacent precisely when the corresponding
edges of $G$ share a common endpoint. Then the number of edges of $L(G)$ equals
\[
\sum_{v} \binom{\deg_G(v)}{2},
\]
where the sum runs over all vertices $v$ of $G$ and $\deg_G(v)$ is the degree of $v$.
\end{theorem}

\begin{theorem}[Paths of length up to the minimum degree]
\lean{SimpleGraph.exists_path_length_of_minDegree}
\label{SimpleGraph.exists_path_length_of_minDegree}
\leanok
Let $G$ be a finite simple graph with minimum degree $\delta(G)$, and let $k$ be a
natural number with $k \le \delta(G)$. Then $G$ contains a path of length exactly $k$;
that is, there exist vertices $u$ and $v$ and a path from $u$ to $v$ in $G$ traversing
exactly $k$ edges.
\end{theorem}

\begin{theorem}[Connectivity via edges crossing every vertex bipartition]
\lean{SimpleGraph.connected_iff_forall_partition_edge}
\label{SimpleGraph.connected_iff_forall_partition_edge}
\leanok
Let $G$ be a simple graph on a finite, nonempty vertex set $V$. Then $G$ is connected if
and only if, for every set $S \subseteq V$ such that both $S$ and its complement $V
\setminus S$ are nonempty, there exist vertices $u \in S$ and $v \in V \setminus S$ that
are adjacent in $G$.
\end{theorem}

\begin{theorem}[Edge count forcing connectedness]
\lean{SimpleGraph.connected_of_edgeCard_gt}
\label{SimpleGraph.connected_of_edgeCard_gt}
\leanok
Let $G$ be a finite simple graph on a vertex set of size $n$, and write $\abs{E(G)}$ for
its number of edges. If
\[
\binom{n-1}{2} < \abs{E(G)},
\]
then $G$ is connected.
\end{theorem}

\begin{theorem}[Complement of a disconnected graph is connected]
\lean{SimpleGraph.compl_connected_of_not_connected}
\label{SimpleGraph.compl_connected_of_not_connected}
\leanok
Let $G$ be a simple graph on a finite, nonempty vertex set. If $G$ is not connected,
then its complement $G^{c}$ — the simple graph on the same vertices in which two
distinct vertices are adjacent precisely when they are non-adjacent in $G$ — is
connected.
\end{theorem}

\begin{theorem}[Adding edges cannot increase the number of connected components]
\lean{SimpleGraph.card_connectedComponent_le_of_le}
\label{SimpleGraph.card_connectedComponent_le_of_le}
\leanok
Let $H$ and $K$ be two simple graphs on a common finite vertex set $V$, and suppose $H
\le K$, meaning that every edge of $H$ is also an edge of $K$. Then $K$ has no more
connected components than $H$; that is, the number of connected components of $K$ is at
most the number of connected components of $H$.
\end{theorem}

\begin{theorem}[Change in component count from deleting an edge]
\lean{SimpleGraph.components_deleteEdge}
\label{SimpleGraph.components_deleteEdge}
\leanok
Let $G$ be a finite simple graph and let $e$ be an edge of $G$. Writing $c(H)$ for the
number of connected components of a graph $H$, and $G - e$ for the graph obtained from
$G$ by removing the edge $e$ (keeping all vertices), one has
\[
c(G) \;\le\; c(G - e) \;\le\; c(G) + 1 .
\]
That is, deleting a single edge either leaves the number of connected components
unchanged or increases it by exactly one.
\end{theorem}

\begin{theorem}[Complement of a large-diameter graph has small diameter]
\lean{SimpleGraph.compl_diam_lt_of_diam_gt}
\label{SimpleGraph.compl_diam_lt_of_diam_gt}
\leanok
Let $G$ be a simple graph on a finite vertex set, and let $G^{\mathrm{c}}$ denote its
complement, the graph on the same vertices in which two distinct vertices are adjacent
exactly when they are non-adjacent in $G$. If the diameter of $G$ is greater than $3$,
then the diameter of $G^{\mathrm{c}}$ is less than $3$.
\end{theorem}

\begin{theorem}[Induced path in a connected non-complete graph]
\lean{SimpleGraph.exists_induced_path_of_connected_not_complete}
\label{SimpleGraph.exists_induced_path_of_connected_not_complete}
\leanok
Let $G$ be a finite simple graph that is connected but not complete. Then there exist
vertices $u$, $v$, $w$ of $G$ such that $u$ is adjacent to $v$ and $v$ is adjacent to
$w$, yet $u$ is not adjacent to $w$.
\end{theorem}

\begin{theorem}[Bipartite graphs are exactly the graphs with no odd cycle]
\lean{SimpleGraph.bipartite_iff_no_odd_cycle}
\label{SimpleGraph.bipartite_iff_no_odd_cycle}
\leanok
Let $G$ be a finite simple graph. Then $G$ admits a proper vertex $2$-colouring —
equivalently, its vertices can be partitioned into two independent sets — if and only if
every cycle in $G$ has even length.
\end{theorem}

\begin{theorem}[Every edge of a closed trail lies on a cycle]
\lean{SimpleGraph.edge_in_cycle_of_closed_trail}
\label{SimpleGraph.edge_in_cycle_of_closed_trail}
\leanok
Let $G$ be a finite simple graph, and let $c$ be a closed trail in $G$ based at a vertex
$u$ — that is, a walk from $u$ to $u$ that traverses no edge more than once — of
positive length. Then every edge $e$ appearing in $c$ lies on a cycle of $G$: there is a
vertex $v$ and a cycle based at $v$ whose edge set contains $e$.
\end{theorem}

\begin{theorem}[Minimum degree at least two forces a cycle]
\lean{SimpleGraph.exists_cycle_of_minDegree_ge_two}
\label{SimpleGraph.exists_cycle_of_minDegree_ge_two}
\leanok
Let $G$ be a finite simple graph whose minimum degree is at least $2$, meaning every
vertex has at least two neighbours. Then $G$ contains a cycle: there is a vertex $v$ and
a closed walk from $v$ to itself that is a cycle.
\end{theorem}

\begin{theorem}[Long cycle from minimum degree]
\lean{SimpleGraph.exists_long_cycle_of_minDegree}
\label{SimpleGraph.exists_long_cycle_of_minDegree}
\leanok
Let $G$ be a finite simple graph whose minimum degree $\delta(G)$ is at least $2$. Then
$G$ contains a cycle of length at least $\delta(G) + 1$.
\end{theorem}

\begin{theorem}[Vertex lower bound for regular graphs of girth four]
\lean{SimpleGraph.girth_four_regular_card_ge}
\label{SimpleGraph.girth_four_regular_card_ge}
\leanok
Let $G$ be a finite simple graph that is $k$-regular, meaning every vertex has exactly
$k$ neighbors, and suppose $G$ has girth $4$, that is, the length of its shortest cycle
is $4$. Then $G$ has at least $2k$ vertices.
\end{theorem}

\begin{theorem}[Moore bound for regular graphs of girth five]
\lean{SimpleGraph.girth_five_regular_card_ge}
\label{SimpleGraph.girth_five_regular_card_ge}
\leanok
Let $G$ be a finite simple graph that is $k$-regular, meaning every vertex has exactly
$k$ neighbours, and suppose $G$ has girth $5$, so that its shortest cycle has length
exactly $5$. Then $G$ has at least $k^2 + 1$ vertices.
\end{theorem}

\begin{theorem}[Cycles in graphs with at least as many edges as vertices]
\lean{SimpleGraph.exists_cycle_of_edgeCard_ge}
\label{SimpleGraph.exists_cycle_of_edgeCard_ge}
\leanok
Let $G$ be a finite simple graph on a nonempty vertex set $V$, and write $n$ for the
number of vertices and $m$ for the number of edges of $G$. If $n \le m$, then $G$
contains a cycle: there exist a vertex $v$ and a closed walk from $v$ to $v$ that is a
cycle.
\end{theorem}
